\begin{textsample}{POS Dim 3 – mg – Score 15.00 – mg/mg\_inf\_camp\_exp\_3.txt}  \label{ex:f3_pos_269}
2.6.1 Campo de Experiências : o Eu, o Outro e o Nós É na interação com os \textbf{pares} e com \textbf{adultos} que as \textbf{crianças} vão constituindo um modo próprio de agir, sentir e pensar e vão descobrindo que existem outros modos de vida, pessoas diferentes, com outros pontos de vista.
Conforme vivem suas primeiras experiências sociais ( na família, na \textbf{instituição} escolar, na coletividade ), constroem \textbf{percepções} e questionamentos sobre si e sobre os outros, diferenciando -se e, simultaneamente, identificando -se como seres individuais e sociais.
Ao mesmo tempo que participam de relações sociais e de \textbf{cuidados} pessoais, as \textbf{crianças} constroem sua autonomia e senso de autocuidado, de reciprocidade e de interdependência com o meio.
Por sua vez, na Educação \textbf{Infantil}, é preciso criar oportunidades para que as \textbf{crianças} entrem em contato com outros grupos sociais e culturais, outros modos de vida, diferentes atitudes, técnicas e rituais de \textbf{cuidados} pessoais e do grupo, costumes, celebrações e narrativas.
Nessas experiências, elas podem ampliar o modo de perceber a si mesmas e ao outro, valorizar sua identidade, respeitar os outros e reconhecer as diferenças que nos constituem como \textbf{seres} humanos. ( BRASIL, 2017, p. 38 ).
A construção da identidade e da autonomia refere -se ao progressivo conhecimento que as \textbf{crianças} vão adquirindo de si mesmas, à autoimagem que se configura por meio deste conhecimento e à capacidade para utilizar recursos pessoais que vão sendo construídos e ampliados, transformando -a e modificando o mundo que a cerca.
De acordo com o Referencial Curricular Nacional para a Educação \textbf{Infantil} ( BRASIL, 1998 ), o termo identidade remete à ideia de distinção.
A identidade é uma \textbf{marca} de diferença entre as pessoas, a começar pelo nome, seguido de todas as características físicas, de modos de agir, de pensar e da história pessoal.
Para Faria e Salles ( 2012 ), a construção da identidade da \textbf{criança} acontece por meio da interação com os outros indivíduos, em diferentes contextos sociais : > [... ] é na relação com o outro, que o sujeito se constitui como ser individual, nas práticas sociais das quais participa, na cultura em que está inserido, apropriando -se dela e transformando -a. [... ] experiências importantes para as \textbf{crianças} relativas à afetividade, ao autoconhecimento, ao \textbf{cuidado} e ao autocuidado, à organização e auto-organização parecem fazer parte do “ currículo oculto ” das \textbf{instituições}, e, desta forma, não têm sido trabalhadas intencionalmente ( FARIA ; SALLES, 2012, p. 101 ).
Construir a identidade implica na ideia de a \textbf{criança} conhecer seus gostos e preferências, habilidades e limites, cultura, sociedade, grupos de convívio e ambientes.
Assim que nasce, o \textbf{bebê} não é capaz de perceber os próprios limites e os limites do outro, o que tornam fundamentais as experiências pelas quais deverá passar para conceber -se como um ser único em meio a outros seres, construindo uma identidade própria.
As situações educativas que a \textbf{criança} vive na \textbf{instituição} e a maneira como são conduzidas pelos professores são muito importantes na formação dos conceitos sobre si mesma.
Os professores, bem como a \textbf{instituição}, devem proporcionar um ambiente rico em interações, que respeite as particularidades de cada indivíduo, reconheça as diversidades e contribua para a construção da unidade coletiva, que favoreça a estruturação da identidade, a construção da sua autonomia e de uma autoimagem positiva.
A autonomia é"a capacidade de se conduzir e de tomar decisões por si próprio, levando, em conta regras, valores, a perspectiva pessoal, bem como a perspectiva do outro". ( BRASIL, 1998 ).
Envolve ações de autocuidado e de decisões.
A \textbf{criança} apresenta -se como um ser com \textbf{vontade} própria e competência para realizar as mais diversas ações.
Assim como o processo de construção da identidade, o desenvolvimento da autonomia tem início desde o \textbf{nascimento} e evolui ao longo da vida.
As \textbf{instituições} de Educação \textbf{Infantil} são espaços privilegiados para as \textbf{crianças} aprenderem e se desenvolver como pessoas e, nesse sentido, torna -se necessário um trabalho intencional com as experiências relacionadas aos saberes e aos conhecimentos sobre si, sobre o outro e sobre o nós.
O docente deve promover em seu planejamento \textbf{diário} situações educativas que potencializam o desenvolvimento da identidade e da autonomia, privilegiando as \textbf{brincadeiras} e as interações como elementos de aprendizagem e de desenvolvimento.
Quando ingressa na \textbf{instituição} de Educação \textbf{Infantil}, a \textbf{criança} encontra possibilidades de ampliar seu olhar e suas experiências, de construir novas formas de relação e de contato com diferentes práticas culturais, criar um repertório de conhecimentos comuns ao grupo, dentre outras conquistas.
O contato das \textbf{crianças} com as \textbf{brincadeiras}, músicas, histórias, jogos e danças, dentre outros, favorece o conhecimento, a valorização da cultura de seu grupo e propicia reflexões sobre a diversidade de \textbf{hábitos}, modos de vida e costumes de seu universo mais próximo e de diferentes épocas, possibilitando, a construção de conhecimentos fundamentados no respeito às diferenças.

% matched lemmas: adulto, bebê, brincadeira, criança, cuidado, diário, hábito, infantil, instituição, marca, nascimento, par, percepção, seres, vontade
\end{textsample}
